\section{Hecke 对应的代数定义}
\label{sec:ch07-hecke-correspondences}
% ============================================================================

第\ref{ch:hecke}章里，Hecke 算子最初是通过双陪集或 $q$-展开公式出现的。
那种定义对计算非常好用：
它能直接告诉我们 Fourier 系数如何变化，
也能快速推出 Euler 乘积与特征形式分解。
但一旦我们开始把模曲线当作 $\Q$ 上的代数曲线，
一个新的障碍立刻出现了：
双陪集公式是解析的、线性的，
而我们现在真正想操纵的是曲线、除子、Jacobian 以及它们在模 $p$ 下的约化。

\textbf{我们遇到了什么问题？}
如果 Hecke 算子只存在于模形式空间里，
那么它们就只是“某些矩阵”。
这样做虽然足够讨论特征值，
却不足以解释这些算子为何与 Jacobian 上的自同态、
与 Galois 作用、
乃至与第\ref{ch:eichler-shimura-l}章将出现的 Frobenius 有关。
所以我们必须找到一个完全几何的定义，
把 $T_p$ 直接写成模曲线上的代数操作。

\textbf{核心思想是什么？}
核心思想是\textbf{对应（correspondence）}\index{对应!algebraic correspondence}。
对一个点 $(E,C_N)$ 来说，
Hecke 算子 $T_p$ 所做的事并不是神秘的线性代数，
而是：
枚举 $E$ 上所有阶 $p$ 子群 $D$，
然后把 $(E,C_N)$ 送到各个商对象
\[
  \bigl(E/D,(C_N+D)/D\bigr).
\]
把这件事同时对所有点执行，
就得到一条位于 $X_0(N)\times X_0(N)$ 内的代数子曲线；
它在两个坐标方向上的投影，就是 Hecke 对应的两支箭头。

\textbf{引入之后能得到什么？}
一旦 Hecke 算子被重新理解为代数对应，
它就自动能够作用在除子、Picard 群与 Jacobian 上；
并且因为整个构造只用到椭圆曲线与有限子群，
它还能在模 $p$ 约化后继续存在。
这正是 Eichler--Shimura 关系的几何起点：
Hecke 与 Frobenius 之所以能放在一起比较，
是因为两者都已经被放进了同一个代数世界里。

\begin{figure}[ht]
\centering
\begin{tikzpicture}[>=Stealth, font=\small, node distance=2.5cm]
  \node[draw, rounded corners, fill=blue!8, minimum width=3.6cm, minimum height=1cm] (x1) {$X_0(N)$};
  \node[draw, rounded corners, fill=green!8, right=of x1, minimum width=4.4cm, minimum height=1cm] (z) {$Z_p=X_0(N;p)$};
  \node[draw, rounded corners, fill=blue!8, right=of z, minimum width=3.6cm, minimum height=1cm] (x2) {$X_0(N)$};

  \draw[->, thick] (z) -- node[above] {$\pi_1(E,C,D)=(E,C)$} (x1);
  \draw[->, thick] (z) -- node[above] {$\pi_2(E,C,D)=\bigl(E/D,(C+D)/D\bigr)$} (x2);

  \node[draw, rounded corners, fill=orange!10, below=1.8cm of z, minimum width=7.0cm, minimum height=0.95cm] (hecke)
    {$T_p=(\pi_2)_*\pi_1^*$，每个点有 $p+1$ 个阶 $p$ 子群};
  \draw[->, thick] (x1.south) -- ++(0,-0.7) -| (hecke.west);
  \draw[->, thick] (x2.south) -- ++(0,-0.7) -| (hecke.east);
\end{tikzpicture}
\caption{Hecke 算子不是凭空写下的矩阵，而是模曲线上的代数对应：先忘掉阶 $p$ 子群，再取商，最后在除子类与 Jacobian 上得到自同态。}
\label{fig:ch07-hecke-correspondences}
\end{figure}


为了突出主要思想，本节先只讨论最常用的情形 $p\nmid N$。

\begin{definition}[Hecke 对应 $T_p$]
\label{def:ch07-hecke-correspondence}
设 $p$ 是不整除 $N$ 的素数。
记 $Z_p$ 为参数化三元组 $(E,C_N,D)$ 的模曲线，
其中 $E$ 是椭圆曲线，$C_N\subset E$ 是阶 $N$ 的循环子群，
$D\subset E$ 是阶 $p$ 的子群。
定义两个态射
\[
  \pi_1(E,C_N,D)=(E,C_N),
  \qquad
  \pi_2(E,C_N,D)=\bigl(E/D,(C_N+D)/D\bigr).
\]
则 $X_0(N)\xleftarrow{\pi_1}Z_p\xrightarrow{\pi_2}X_0(N)$
称为\textbf{Hecke 对应 $T_p$}。
\end{definition}

\begin{proposition}[两个投影的次数都是 $p+1$]
\label{prop:ch07-hecke-degree-p-plus-1}
在 \cref{def:ch07-hecke-correspondence} 的情形下，
$\pi_1$ 与 $\pi_2$ 都是有限态射，且
\[
  \deg \pi_1=\deg \pi_2=p+1.
\]
\end{proposition}

\begin{proof}
先看 $\pi_1$。
固定 $X_0(N)$ 上的一个非尖点 $(E,C_N)$。
因为 $p\nmid N$，
$E[p]\cong (\Z/p\Z)^2$，
其中阶 $p$ 子群恰好对应于二维向量空间中的一维子空间。
一维子空间的个数是
\[
  \frac{p^2-1}{p-1}=p+1.
\]
因此 $\pi_1^{-1}(E,C_N)$ 恰有 $p+1$ 个点，
分别对应于 $E$ 的 $p+1$ 个阶 $p$ 子群 $D$。
这说明一般纤维大小为 $p+1$，故 $\deg\pi_1=p+1$。
有限性则由纤维有限且模曲线射影可知。

再看 $\pi_2$。
固定目标点 $(E',C_N')$。
要描述其原像，等价于枚举所有次数为 $p$ 的循环同源
\[
  \varphi:E\to E'
\]
以及源端的阶 $N$ 循环子群 $C_N$，满足 $\varphi(C_N)=C_N'$。
由对偶同源理论，
每条这样的同源都对应于 $E'$ 上一个阶 $p$ 子群 $D'\subset E'$：
取对偶同源 $\widehat\varphi:E'\to E$，其核就是 $D'$。
反过来，从任意阶 $p$ 子群 $D'\subset E'$ 出发，
取商映射 $E'\to E'/D'$ 并再取对偶，
便恢复一条次数 $p$ 的同源回到 $E'$。
因此 $\pi_2$ 的一般纤维也由 $E'$ 的阶 $p$ 子群控制，仍有 $p+1$ 个。
故 $\deg\pi_2=p+1$。
\end{proof}

\begin{definition}[对应诱导的除子映射]
\label{def:ch07-correspondence-on-divisors}
设一般地有两条光滑射影曲线之间的对应
\[
  X\xleftarrow{\alpha} Z\xrightarrow{\beta} Y,
\]
其中 $\alpha,\beta$ 都是有限态射。
定义除子上的群同态
\[
  \beta_*\alpha^*:\Div(X)\longrightarrow \Div(Y)
\]
为先拉回后推出。
当 $X=Y=X_0(N)$ 且 $Z=Z_p$ 时，
这就是代数意义下的 Hecke 算子 $T_p$。
\end{definition}

\begin{proposition}[Hecke 对应作用在 Jacobian 上]
\label{prop:ch07-hecke-on-jacobian}
\cref{def:ch07-hecke-correspondence} 给出的映射
\[
  T_p=(\pi_2)_*\pi_1^*:
  \Div\bigl(X_0(N)\bigr)\longrightarrow \Div\bigl(X_0(N)\bigr)
\]
保持主除子与次数零除子，
故诱导 Jacobian 上的自同态
\[
  T_p:J_0(N)\longrightarrow J_0(N).
\]
而且这个自同态定义在 $\Q$ 上。
\end{proposition}

\begin{proof}
先看次数。
若 $D=\sum n_P P$ 是 $X_0(N)$ 上的除子，
则由有限态射的定义可知
\[
  \deg \pi_1^*D=(\deg\pi_1)\deg D.
\]
特别地，当 $\deg D=0$ 时，拉回后仍次数为零。
另一方面，
若 $E=\sum m_Q Q$ 是 $Z_p$ 上的除子，
推出只是把同像点上的系数相加，
故
\[
  \deg (\pi_2)_*E=\deg E.
\]
于是 $T_p$ 保持次数零除子。

再看主除子。
若 $D=\divisor(f)$，则
\[
  \pi_1^*D=\divisor(f\circ \pi_1)
\]
仍是主除子。
对 $Z_p$ 上任意有理函数 $g$，
有限态射 $\pi_2$ 的范数 $\Nm_{\pi_2}(g)$ 是 $X_0(N)$ 上的有理函数，
并满足
\[
  (\pi_2)_*\divisor(g)=\divisor\bigl(\Nm_{\pi_2}(g)\bigr).
\]
因此推出也把主除子送到主除子。

综上，$T_p$ 在 $\operatorname{Pic}^0(X_0(N))\cong J_0(N)$ 上良定义。
又因为 $Z_p,\pi_1,\pi_2$ 都是由模解释定义的代数对象，
它们对 Galois 共轭稳定，
故诱导的 Jacobian 自同态也定义在 $\Q$ 上。
\end{proof}

\begin{proposition}[代数定义与经典 Hecke 算子一致]
\label{prop:ch07-algebraic-equals-classical-hecke}
在识别
\[
  H^0\bigl(X_0(N),\Omega^1\bigr)\cong \Sk_2\bigl(\GammaO(N)\bigr)
\]
下，\cref{prop:ch07-hecke-on-jacobian} 中的代数算子 $T_p$
在微分上的作用，
与第\ref{ch:hecke}章定义的经典 Hecke 算子相同。
\end{proposition}

\begin{proof}
对一般对应
\[
  X\xleftarrow{\alpha} Z\xrightarrow{\beta} X,
\]
其在 Jacobian 上的作用是 $(\beta)_*\alpha^*$。
在原点余切空间上线性化，
得到的正是微分形式上的复合映射
\[
  H^0(X,\Omega^1)\xrightarrow{\ \alpha^*\ } H^0(Z,\Omega^1)
  \xrightarrow{\ \beta_*\ } H^0(X,\Omega^1),
\]
其中 $\beta_*$ 是迹映射。

当 $X=X_0(N)$、$Z=Z_p$ 时，
$\alpha=\pi_1$ 先把一个点 $(E,C_N)$ 拉回到所有可能的三元组 $(E,C_N,D)$；
接着 $\beta=\pi_2$ 把它们分别推到各个商对象
\[
  \bigl(E/D,(C_N+D)/D\bigr).
\]
这恰好就是第\ref{ch:hecke}章“枚举次数 $p$ 同源”的定义。
因此在 $H^0(X_0(N),\Omega^1)$ 上得到的作用与经典 $T_p$ 完全一致。
再由 \cref{thm:s2-holomorphic-differentials} 的同构，
这就是 $\Sk_2(\GammaO(N))$ 上的 Hecke 算子。
\end{proof}

\textbf{注。}
这一节的最大收获，不是重新得到一个旧定义，
而是把 Hecke 算子放回到了几何世界里。
从现在起，$T_p$ 不只是“对 $q$-展开做某种线性变换”，
而是模曲线上的代数对应。
因此它既能作用在 Jacobian 上，
也能在整模型与特殊纤维上继续存在。
特别是在 $p\nmid N$ 时，
第\ref{ch:eichler-shimura-l}章会看到：
约化到 $\F_p$ 后，Hecke 对应与 Frobenius 之间会出现惊人的精确关系。

\textbf{模 $p$ 约化视角。}
当 $p\nmid N$ 时，$X_0(N)$ 在 $\F_p$ 上仍然是光滑曲线，
因此 Hecke 对应可以逐项约化到特殊纤维上。
在特征 $p$ 中，阶 $p$ 子群不再只是 $p+1$ 个离散点的集合，
而会出现 Frobenius 与 Verschiebung 这两条最基本的次数 $p$ 同源。

\begin{proposition}[特殊纤维上的 $T_p$]
\label{prop:hecke-correspondence-mod-p}
设 $p\nmid N$。
则 Hecke 对应 $T_p$ 在 $X_0(N)_{\F_p}$ 上仍有代数意义：
对一个几何点 $(E,C_N)$，它枚举 $E$ 上所有阶 $p$ 子群概形 $D$，
并把 $(E,C_N)$ 送到各个商对象
\[
  \bigl(E/D,(C_N+D)/D\bigr).
\]
在普通 locus 上，这个对应分裂为两条最重要的分支：
一条是相对 Frobenius 的图像，另一条是 Verschiebung 的图像。
\end{proposition}

\begin{proof}[说明]
由第7章整模型的存在性，当 $p\nmid N$ 时模曲线与对应都能约化到 $\F_p$ 上。
于是 $T_p$ 的定义仍然是“枚举阶 $p$ 子群再取商”。
在特征 $p$ 中，这些子群里有两个最几何的方向：
一个是 $\ker(\Frob_p)$，商出去得到相对 Frobenius；
另一个对应于 Verschiebung 的对偶核。
因此在特殊纤维上，Hecke 对应恰好显现为 Frobenius 与 Verschiebung 的叠加。
\end{proof}

\begin{corollary}[代数版 Eichler--Shimura 关系]
\label{cor:algebraic-eichler-shimura-correspondence}
在 $X_0(N)_{\F_p}$ 的对应代数里，
有
\[
  T_p=\Frob_p+V_p,
\]
其中 $V_p$ 表示 Verschiebung 对应。
把 $V_p$ 写成 $\langle p\rangle\Frob_p^{-1}$ 的记号后，
这就是常写的代数关系
\[
  T_p\equiv \Frob_p+\langle p\rangle\Frob_p^{-1}.
\]
它正是第8章 Eichler--Shimura 关系的代数几何起点。
\end{corollary}

\begin{proof}[说明]
上一命题已经指出：在特征 $p$ 的模解释里，
$T_p$ 的两条基本分支就是 Frobenius 与 Verschiebung。
而在对应代数中，Verschiebung 与 Frobenius 互为“反向”的次数 $p$ 对应；
在带级别结构的情形，需要再记住菱形算子 $\langle p\rangle$ 对循环子群的作用，
于是常把第二支写成 $\langle p\rangle\Frob_p^{-1}$。
在 Jacobian 或 $\ell$-进上同调上，这进一步转化为
\[
  X^2-T_pX+p\langle p\rangle=0,
\]
也就是下一章将真正使用的 Eichler--Shimura 关系。
\end{proof}

\textbf{注。}
这里的 $\Frob_p^{-1}$ 是对应代数中的逆向记号，
不应理解为曲线态射意义下的真反函数。
对 $\Gamma_0(N)$ 的平凡 Nebentypus 部分，$\langle p\rangle$ 最终按标量 $1$ 作用，
所以在权 $2$ 的经典表述里，常把这条关系简写成
\[
  T_p=\Frob_p+p\Frob_p^{-1}.
\]
